\section{Fourier amplification of partial permutation laws}
\label{shear:sec:completion}

The marginal estimates now provide two routes to coset information.
A bound for every fixed input list applies to each prescribed partition
whose complementary lists meet the stated length condition. An averaged
input bound first requires the choice of one partition. In either case
the Fourier argument uses one common
modification of the permutation law whose relevant coset masses are
bounded. We first state that modification and the transfer estimates,
then apply them to the preceding marginal bounds.

\subsection{From complementary marginals to one retained law}

Let \(X\) be a finite label set, \(G=\Sym(X)\), and let
\(M_1,\ldots,M_r\) be a partition of \(X\) into nonempty parts. Write
\(H_i=\Sym(M_i)\) for the subgroup fixing \(X\setminus M_i\) pointwise.
Two label-to-position permutations agree on \(X\setminus M_i\) exactly
when they lie in the same left coset \(gH_i\). Thus the complementary
image tuple identifies that coset, and its distance from a uniform
injection is the distance of the coset law from uniform on \(G/H_i\).
Write \(\delta_i\) for this distance and put
\(\delta=\sum_i\delta_i\).

If a marginal estimate averages over omitted sets of size \(|X|/r\),
sample a uniform ordered partition into \(r\) equal parts. Each part
has the required uniform subset law, so the expected sum of its
complementary errors is \(r\) times the mean single-subset error.
Some deterministic partition has a sum no larger than this expectation.
The ordering chosen on a complementary set does not affect its error,
because changing that ordering permutes the tuple coordinates.
Thus an average over uniform ordered input tuples is also the
corresponding average over complementary sets.

For the representation estimates, write \(\rho_\lambda\) for the
symmetric-group irreducible indexed by \(\lambda\) and
\(D_\lambda=f^\lambda=\dim\rho_\lambda\) for its degree. For \(u>0\), define
\[
 C_u=\sup_{s\ge1}\sum_{\lambda\vdash s}(f^\lambda)^{-u},
 \qquad
 Z_u(s)=\sum_{\substack{\lambda\vdash s\\
                   \lambda\notin\{(s),(1^s)\}}}(f^\lambda)^{-u}.
\]
The exclusions form a set, so the sole partition of \(1\) is removed once.

\begin{proposition}[Degree and coset inputs]
\label{shear:prop:fourier-inputs}
For every fixed \(u>0\), \(C_u<\infty\) and \(Z_u(s)\to0\) as \(s\to\infty\).

Let \(H_1,\ldots,H_r\) be subgroups of a finite group \(G\), and let
\(\mu\) be a probability measure. Write \(\mu_H\) for its image on
\(G/H\). If
\[
 \delta=\sum_{i=1}^r
 \|\mu_{H_i}-U_{G/H_i}\|_{\TV},
\]
there is a measure \(0\le f\le\mu\), of mass \(z\ge1-\delta\), with
\[
 f(gH_i)\le [G:H_i]^{-1}\qquad(g\in G,\ 1\le i\le r).
\]
It can be obtained by successive trimming of excess coset mass or by
taking the pointwise minimum of the separately trimmed measures.
When \(z>0\), its normalization is within \(1-z\) of \(\mu\) in TV and
has coset cap \(z^{-1}\) times uniform.

Alternatively, delete the union of all cosets whose original \(\mu\)-mass
exceeds twice uniform. The deleted mass is \(\eta\le2\delta\).
If \(\delta<1/2\), the normalized retained law is exactly \(\eta\) away
from \(\mu\) in TV and has coset cap \(2/(1-\eta)\) times uniform.
\end{proposition}
\begin{proof}
The degree assertions follow from
\cite[Theorem~\Lref{l:degree-all-powers}]{partial-fourier}.
The same companion gives separate proofs at powers \(1,2,10\) in
Lemmas~\Lref{l:degree-inverse-first},
\Lref{l:degree-inverse-square}, and \Lref{l:degree-inverse-ten};
these also justify the fixed-power sums used below.
The two modifications, including their common-measure and mass claims,
follow from \cite[Lemma~\Lref{l:trim}]{partial-fourier}.
The normalization assertion for \(f\) follows from
Lemma~\ref{lem:subprobability-fourier} below; dividing the exact cap
by \(z\) gives the normalized cap.
\end{proof}

\subsection{Fourier bounds for disjoint supports}

For a nonnegative mass function \(f\) on \(G\), use
\[
 \widehat f(\rho)=\sum_{g\in G}f(g)\rho(g).
\]
The Hilbert--Schmidt norm is unnormalized:
\(\|A\|_{\HS}^2=\tr(A^*A)\).
The next proposition extends the companion's partition estimates to
disjoint supports that may leave labels unused. The common restriction
decomposition has two uses: assigning each product type to one small
factor controls the orbit of a vector, while overlapping central
projections control the whole Fourier matrix.

\begin{proposition}[Fourier bounds from disjoint coset systems]
\label{shear:prop:disjoint-transfer}
Let \(X\) have \(N\) elements and \(G=\Sym(X)\). Let
\(M_1,\ldots,M_r\) be disjoint nonempty subsets of \(X\), where \(r\ge1\),
and put \(H_i=\Sym(M_i)\), fixing the complement pointwise.
Their union need not be \(X\). Suppose \(f\ge0\) has mass \(z\le1\) and
\[
 f(gH_i)\le K[G:H_i]^{-1}\qquad(g\in G,\ 1\le i\le r)
\]
for some \(K\ge0\). For an irreducible unitary representation \(\rho\)
of degree \(D\), set
\[
 S_i(D)=
 \sum_{\substack{\tau\in\widehat H_i\\ \dim\tau\le D^{1/r}}}
       (\dim\tau)^2.
\]
For every \(u>0\), the separate orbit and isotypic estimates are
\begin{align}
 \|\widehat f(\rho)\|_{\op}
 &\le \left(\frac{zK}{D}\sum_iS_i(D)\right)^{1/2}
 \le \sqrt{zKrC_u}\,D^{-1/2+(u+2)/(2r)},
 \label{shear:eq:noncovering-orbit-general}\\
 \|\widehat f(\rho)\|_{\HS}^2
 &\le \frac{zK}{D}\sum_iS_i(D)
 \le zKrC_uD^{-1+(u+2)/r}.
 \label{shear:eq:noncovering-isotypic-general}
\end{align}
In particular, when \(r>4\), the inverse-square orbit estimate gives
\begin{equation}
 \|\widehat f(\rho)\|_{\op}
 \le\sqrt{rKC_2}\,D^{-(1-4/r)/2}.
 \label{shear:eq:orbit-bound}
\end{equation}
A separate coefficient-projection argument gives
\begin{equation}
 \|\widehat f(\rho)\|_{\HS}^2
 \le\frac{K^2}{D}\sum_{i=1}^r
       \sum_{\substack{\tau\in\widehat H_i\\
                            \dim\tau\le D^{1/r}}}(\dim\tau)^2
 \le K^2rC_{10}D^{-1+12/r}.
 \label{shear:eq:hs-general}
\end{equation}
Finally, for \(\rho=\rho_\lambda\) with \(\lambda\vdash N\), put
\[
 k_\lambda=N-\max(\lambda_1,\lambda'_1),
 \qquad M(k)=\sum_{j=0}^k p(j),
\]
where \(p(j)\) is the partition number and \(p(0)=1\). The branching
refinement is
\begin{equation}
 \|\widehat f(\rho_\lambda)\|_{\HS}^2
 \le zKrM(k_\lambda)D^{-1+2/r}.
 \label{shear:eq:hs-branching}
\end{equation}
All these statements allow unequal support sizes and unrestricted
restriction multiplicities. In particular they include the
probability-law cases \(z=1\).
\end{proposition}

\begin{proof}[Proof of Proposition~\ref{shear:prop:disjoint-transfer}]
Write \(A=\widehat f(\rho)\) and abbreviate \(S_i(D)\) to \(S_i\).
If \(z=0\), then \(f=0\) and every assertion is immediate. Assume \(z>0\).

\paragraph{The common small-type cover.}
Disjoint supports embed \(H_1\times\cdots\times H_r\) in \(G\):
the factors commute, and a product that is identity restricts to the
identity on every \(M_i\). Complete reducibility and the
direct-product description of irreducibles
\cite[Theorems~3.12 and~4.25 and Proposition~3.14]{etingof} give
\[
 V_\rho\big|_{H_1\times\cdots\times H_r}
 =
 \bigoplus_{\boldsymbol\tau}
 (V_{\tau_1}\otimes\cdots\otimes V_{\tau_r})
       \otimes\mathcal A_{\boldsymbol\tau}.
\]
The sum is over distinct occurring product types, with their full
multiplicity spaces. Each has
\(\prod_i\dim\tau_i\le D\), so some factor has degree at most
\(D^{1/r}\). Unused labels can change the occurring types and their
multiplicities, but do not change this product bound.

Let \(P_{i,\tau}\) be the orthogonal projection onto all copies of
the \(H_i\)-type \(\tau\), and put
\[
 P_i=\sum_{\substack{\tau\in\widehat H_i\\
                         \dim\tau\le D^{1/r}}}P_{i,\tau},
 \qquad Q_i=P_i\prod_{j<i}(I-P_j).
\]
On each product-type summand, \(P_i\) is identity or zero according
as the \(i\)th factor is small or large. Hence the \(P_i\) commute,
and at least one is identity on each summand. Therefore
\begin{equation}
 I\preceq\sum_{i=1}^rP_i,\qquad
 \sum_{i=1}^rQ_i=I,\qquad Q_iQ_j=0\quad(i\ne j).
 \label{shear:eq:disjoint-projector-cover}
\end{equation}
The \(Q_i\) assign a whole product type to its first small factor;
the \(P_i\) may overlap. Both families retain every multiplicity.
For every \(u>0\), the definition of \(C_u\) gives
\[
 S_i\le D^{(u+2)/r}
       \sum_{\tau\in\widehat H_i}(\dim\tau)^{-u}
 \le C_uD^{(u+2)/r}.
\]

\paragraph{The orbit estimate.}
For \(v\in V_\rho\), write \(v_i=Q_iv\) and
\(W_i=\Span\{\rho(h)v_i:h\in H_i\}\). The single-vector orbit lemma
\cite[Lemma~\Lref{l:single-orbit}]{partial-fourier} gives
\(\dim W_i\le S_i\). Indeed, in all copies of a type \(\tau\), the
orbit of its component is contained in the image of
\[
 \operatorname{End}(V_\tau)\longrightarrow
 V_\tau\otimes\mathcal A_\tau,\qquad
 T\longmapsto(T\otimes I)v_{i,\tau},
\]
whose dimension is at most \((\dim\tau)^2\), independently of the
multiplicity.

Let \(\Pi_i\) project onto \(W_i\). Since \(W_i\) is \(H_i\)-invariant,
the projection onto \(\rho(g)W_i\) depends only on \(gH_i\).
The coset cap and Schur averaging yield
\[
 \sum_g f(g)\|\Pi_{\rho(g)W_i}w\|^2
 \le K\frac{\dim W_i}{D}\|w\|^2.
\]
Here the uniform average of
\(\rho(g)\Pi_i\rho(g)^*\) is \((\dim W_i/D)I\), by its trace and
irreducibility. Weighted Cauchy--Schwarz retains the mass \(z\):
\[
 |\langle w,Av_i\rangle|^2
 \le z\|v_i\|^2\sum_g f(g)\|\Pi_{\rho(g)W_i}w\|^2
 \le\frac{zKS_i}{D}\|w\|^2\|v_i\|^2.
\]
The \(v_i\) are orthogonal. Summing and applying Cauchy--Schwarz over
\(i\) gives
\[
 |\langle w,Av\rangle|
 \le\left(\frac{zK}{D}\sum_iS_i\right)^{1/2}\|w\|\|v\|.
\]
This proves \eqref{shear:eq:noncovering-orbit-general}.
Its specialization \(u=2\), with \(z\le1\), is
\eqref{shear:eq:orbit-bound}; the condition \(r>4\) makes that
dimension exponent negative.

\paragraph{The isotypic estimate.}
For a subgroup \(H\le G\) and an irreducible \(H\)-type \(\tau\) of
degree \(a\), define
\[
 e_{H,\tau}(h)=\frac{a}{|H|}\overline{\chi_\tau(h)}
 \quad(h\in H),\qquad e_{H,\tau}=0\quad\text{off }H.
\]
Its transform is the projection \(P_{\sigma,H,\tau}\) onto all copies
of \(\tau\) in \(\sigma|_H\). Thus
\(\widehat{f*e_{H,\tau}}(\sigma)=
\widehat f(\sigma)P_{\sigma,H,\tau}\): a bound on the left cosets
\(gH\) controls a projection on the right of the Fourier matrix.
The companion's cosetwise isotypic lemma
\cite[Lemma~\Lref{l:central-isotypic}]{partial-fourier} applies to
every subgroup \(H\). If \(f(gH)\le K[G:H]^{-1}\), it gives
\begin{equation}
 \sum_{\sigma\in\widehat G}D_\sigma
       \|\widehat f(\sigma)P_{\sigma,H,\tau}\|_{\HS}^2
 =|G|\sum_g|(f*e_{H,\tau})(g)|^2\le Kza^2.
 \label{shear:eq:noncovering-central}
\end{equation}
This is the central-projection step; it retains the whole projected
Fourier matrix and the mass factor.

Trace the first inequality of \eqref{shear:eq:disjoint-projector-cover}
against \(A^*A\), and apply \eqref{shear:eq:noncovering-central} to
each selected \(H_i\)-type after discarding the other ambient
irreducibles. Then
\[
 D\|A\|_{\HS}^2
 \le D\sum_i\sum_{\dim\tau\le D^{1/r}}
       \tr(A^*AP_{i,\tau})
 =D\sum_i\sum_{\dim\tau\le D^{1/r}}
       \|AP_{i,\tau}\|_{\HS}^2
 \le Kz\sum_iS_i.
\]
The trace inequality uses positivity of
\(A(\sum_iP_i-I)A^*\); it does not require \(A^*A\) to commute with
the projections. This proves
\eqref{shear:eq:noncovering-isotypic-general} using the overlapping
\(P_i\), rather than the orbit assignment \(Q_i\).

\paragraph{The coefficient alternative.}
The independent coefficient calculation in
\cite[Lemma~\Lref{l:isotypic-alternatives}]{partial-fourier}
bounds each coset separately. Its norm and orientation give the
selected-type refinement in \eqref{shear:eq:hs-general}.
For one subgroup \(H\le G\), put \(F=|G|f\), so
\(\E_{h\sim U_H}F(xh)=[G:H]f(xH)\le K\).
For an \(H\)-type \(\tau\) of degree \(a\), choose one representative
\(x\) of each coset and set
\[
 T_x=\E_{h\sim U_H}F(xh)\tau(h),\qquad
 q(xh)=a\sum_{j,\ell=1}^a(T_x)_{j\ell}\overline{\tau(h)_{j\ell}}.
\]
The function \(q\) on this coset is the orthogonal projection of
\(F\) onto the conjugate coefficient space of \(\tau\). Schur
orthogonality and the nonnegative weights \(F(xh)\) give
\begin{equation}
 \E_{h\sim U_H}|q(xh)|^2
 =a\|T_x\|_{\HS}^2
 \le a^2\|T_x\|_{\op}^2
 \le a^2\bigl(\E_{h\sim U_H}F(xh)\bigr)^2
 \le K^2a^2.
 \label{shear:eq:coset-coefficient-norm}
\end{equation}
Expanding \(T_x\) and using, for \(h,k\in H\),
\(\overline{\chi_\tau(k^{-1}h)}=\chi_\tau(kh^{-1})\)
shows that these coset projections form the single function
\(q=|G|(f*e_{H,\tau})\). Thus their tested Fourier matrix is
\(\widehat f(\sigma)P_{\sigma,H,\tau}\), with the same right-hand
orientation as above. Averaging the preceding norm bound over cosets
and applying complex finite-group Plancherel gives
\begin{equation}
 \sum_{\sigma\in\widehat G}D_\sigma
       \|\widehat f(\sigma)P_{\sigma,H,\tau}\|_{\HS}^2
 =\E_{g\sim U_G}|q(g)|^2\le K^2a^2.
 \label{shear:eq:coefficient-one-type}
\end{equation}
The calculation counts a type once, regardless of its multiplicity,
and assumes no pointwise bound inside the coset. The same positive
trace comparison now gives
\(\|A\|_{\HS}^2\le K^2D^{-1}\sum_iS_i\).
Using \(S_i\le C_{10}D^{12/r}\) proves
\eqref{shear:eq:hs-general}, including its original sum over the
types of the unmodified subgroups \(H_i\).

\paragraph{The branching refinement.}
Each \(H_i\) is conjugate to the standard subgroup \(S_{|M_i|}\)
of \(S_N\). Young branching
\cite[Theorem~9.2]{James1978} shows that every occurring type
\(\tau\vdash |M_i|\) is a subdiagram of \(\lambda\); different
corner-removal sequences can give arbitrary multiplicity.
If \(\lambda_1\ge\lambda'_1\), the diagram \(\tau\) has at most
\(k_\lambda\) boxes below its first row. Its lower diagram, together
with its fixed size, determines the first row, so at most
\(M(k_\lambda)\) distinct types can occur. If
\(\lambda'_1>\lambda_1\), transpose both diagrams for this count.
Transposition preserves degrees and does not replace \(f\) or its
Fourier matrix by a sign-twisted law.

Nonoccurring isotypic projections are zero. In the positive trace
comparison, apply \eqref{shear:eq:noncovering-central} only to the
occurring selected types and use
\((\dim\tau)^2\le D^{2/r}\). There are at most \(M(k_\lambda)\) of
them for each support, giving \eqref{shear:eq:hs-branching}.
The same argument includes \(k_\lambda=0\), where the representation
is trivial or sign, \(D=1\), and \(M(0)=1\). It also includes all
unequal nonempty supports and every restriction multiplicity.
\end{proof}


The isotypic estimate has the same dimension exponent as the squared
orbit estimate, but bounds the whole Hilbert--Schmidt norm. The coefficient
estimate records a different proof of a bound with cap factor \(K^2\).
The branching estimate counts only types that can occur in a fixed
ambient diagram. These methods use only coset masses; none requires a
pointwise bound on the conditional distribution inside a coset.

\subsection{Normalization and independent time blocks}

We record the mass normalization once for all the applications.
The Fourier and sign conventions are those of the completed local proof.

\begin{lemma}[Fourier normalization for retained mass]
\label{lem:subprobability-fourier}
For a complex mass function \(a\) on a finite group \(G\),
\begin{equation}
 \sum_\rho D_\rho\|\widehat a(\rho)\|_{\HS}^2
 =|G|\sum_g|a(g)|^2.
 \label{shear:eq:plancherel-mass}
\end{equation}
If \(\xi\ge0\) has mass \(m\le1\), its transform is not divided by \(m\),
and
\begin{align}
 |G|\sum_g\left|\xi(g)-\frac{m}{|G|}\right|^2
 &=\sum_{\rho\ne\mathrm{triv}}D_\rho
                   \|\widehat\xi(\rho)\|_{\HS}^2,
 \label{eq:subprobability-centered}\\
 \left(\sum_g|\xi(g)-U_G(g)|\right)^2
 &\le(1-m)^2+\sum_{\rho\ne\mathrm{triv}}D_\rho
                   \|\widehat\xi(\rho)\|_{\HS}^2.
 \label{eq:subprobability-full}
\end{align}
If \(\xi\le\mu\) for a probability measure \(\mu\), then
\(\sum_g|\mu(g)-\xi(g)|=1-m\). When \(m>0\),
\(\|\xi/m-\mu\|_{\TV}\le1-m\).
\end{lemma}
\begin{proof}
The regular-character expansion in Lemma~\ref{lem:plancherel}, with
coefficients \(a(g)\overline{a(h)}\), proves
\eqref{shear:eq:plancherel-mass}. Apply it to \(\xi-mU_G\) and
\(\xi-U_G\). Their trivial coefficients are respectively \(0\) and
\(m-1\), and their nontrivial coefficients are those of \(\xi\).
Cauchy--Schwarz gives \eqref{eq:subprobability-full}.
If \(\xi\le\mu\), the nonnegative difference has mass \(1-m\).
Also \(\sum_g|\xi(g)/m-\xi(g)|=1-m\); the triangle inequality and
the factor \(1/2\) in TV prove the normalization claim.
\end{proof}

For a subprobability \(f\) of mass \(z\) and a positive integer \(\ell\),
the same identity and the convolution rule give
\begin{equation}
 |G|\sum_g|f^{*\ell}(g)-|G|^{-1}|^2
 =(z^\ell-1)^2+
 \sum_{\rho\ne\mathrm{triv}}D_\rho
             \|\widehat f(\rho)^\ell\|_{\HS}^2.
 \label{shear:eq:subprob-plancherel}
\end{equation}
One half of the square root bounds one half of the \(\ell^1\) distance
from \(f^{*\ell}\) to \(U_G\); for a probability this is TV.
The matrix bounds used below are
\[
 \|A^\ell\|_{\HS}\le\sqrt D\,\|A\|_{\op}^{\ell},
 \qquad
 \|A^\ell\|_{\HS}\le\|A\|_{\HS}^{\ell}.
\]
They require no normality assumption.

For every positive-time physical block law \(\mu\),
Lemma~\ref{lem:fair-sign} gives \(\widehat\mu(\sgn)=0\).
If a probability \(\nu\) is within \(\eta\) of \(\mu\) in TV, then
\begin{equation}
 |\widehat\nu(\sgn)|\le2\eta.
 \label{shear:eq:parity}
\end{equation}
If instead \(0\le f\le\mu\) has mass \(z\), then
\(|\widehat f(\sgn)|\le1-z\). Positivity also gives
\(f^{*\ell}\le\mu^{*\ell}\), with missing mass \(1-z^\ell\).
For nearby probabilities, replacing the \(\ell\) factors one at a time
costs at most \(\ell\eta\) in TV.

Every block length below is a multiple of \(d\). Thus
\(R^T=I\) in \eqref{eq:frames:cyclic}, and the coordinate block
\(\mathcal L(X_T)\) is exactly the physical law \(q_d^{*T}\).
Independent time blocks combine by convolution.

\begin{corollary}[Two time blocks from the isotypic transfer]
\label{shear:cor:512}
For the Thorp shuffle on \(N=2^d\) cards,
\[
 \|q_d^{*(512d)}-U_G\|_{\TV}\longrightarrow0
 \qquad(d\longrightarrow\infty).
\]
The convergence is uniform over deterministic initial orderings.
\end{corollary}
\begin{proof}
Let \(T=256d\) and \(\mu=q_d^{*T}\). For \(d\ge3\), partition the labels
into eight sets \(M_i\) of size \(N/8\). The fixed-list estimate
\eqref{shear:eq:256-marginal} bounds each complementary image-tuple
error by
\[
 \delta_N=\frac{N^{3/2}}2
 \left((31/32)^{255d}+64e^{-N/64}\right)^{1/2}=o(1).
\]
Exact trimming by Proposition~\ref{shear:prop:fourier-inputs} gives
\(0\le f\le\mu\) with mass \(z\), where \(1-z\le8\delta_N=o(1)\), and
cap one on all eight coset systems.

For this covering partition, the companion isotypic estimate
\cite[Proposition~\Lref{l:isotypic}]{partial-fourier}, with inverse-degree
parameter \(1\), gives
\(\|\widehat f(\rho)\|_{\HS}^2\le8zC_1D^{-5/8}\le8C_1D^{-5/8}\).
Consequently
\[
 \sum_{\substack{\lambda\vdash N\\
                  \lambda\notin\{(N),(1^N)\}}}
 D_\lambda\|\widehat f(\rho_\lambda)^2\|_{\HS}^2
 \le(8C_1)^2Z_{1/4}(N)=o(1),
\]
since \(1-2(5/8)=-1/4\) and
\cite[Theorem~\Lref{l:degree-all-powers}]{partial-fourier}
gives the inverse-quarter limit. The sign coefficient of \(f\) is at
most \(1-z\) in magnitude, and the trivial coefficient of \(f^{*2}\)
is \(z^2\). Lemma~\ref{lem:subprobability-fourier} therefore gives
\[
 \left(\sum_g|f^{*2}(g)-U_G(g)|\right)^2
 \le(1-z^2)^2+(1-z)^4+(8C_1)^2Z_{1/4}(N)=o(1).
\]
Adding back the missing mass \(1-z^2=o(1)\) proves
\(\|\mu^{*2}-U_G\|_{\TV}=o(1)\). The two physical blocks take
\(2T=512d\) shuffles, and \eqref{eq:worst-state} gives uniformity over starts.
\end{proof}

The orbit, coefficient, and branching estimates give the following
separate applications, with different modifications of the law.

\begin{theorem}[Eight time blocks from the orbit estimate]
\label{shear:thm:2048}
For the physical shuffle on \(N=2^d\) labeled cards,
\begin{equation}
 \|q_d^{*(2048d)}-U_G\|_{\TV}=o(1).
 \label{shear:eq:2048}
\end{equation}
The fixed-list estimate at \(257d\) also gives a proof using eight time blocks
at \(2056d\) with a retained subprobability. The random omitted set
estimate gives another proof at \(2056d\) using normalized deletion.
All three conclusions hold for every deterministic starting deck.
\end{theorem}
\begin{proof}
For \(T=256d\), let \(\mu=q_d^{*T}\) and use eight equal label parts.
The error \(\delta_N\) in the preceding proof bounds each complementary
marginal. Deleting heavy cosets loses
\(\eta\le16\delta_N=o(1)\) and produces a probability \(\nu\) with
cap \(K\le4\) for large \(d\). Equation~\eqref{shear:eq:orbit-bound}
with \(r=8\) gives
\[
 \|\widehat\nu(\rho)\|_{\op}
 \le C D^{-1/4},\qquad C=\sqrt{32C_2}.
\]
For eight factors the nonsign, nontrivial Plancherel contribution is
at most
\[
 C^{16}\sum_{\substack{\lambda\vdash N\\
                        \lambda\notin\{(N),(1^N)\}}}
 D_\lambda^{\,2-16/4}=C^{16}Z_2(N)=o(1).
\]
The sign contribution is at most \((2\eta)^{16}\), and the trivial
coefficient is one. Thus \(\nu^{*8}\) tends to uniform in TV.
Restoring the true factors costs at most \(8\eta=o(1)\), and
\(8T=2048d\) physical shuffles.

For the fixed-list subprobability route take \(T=257d\).
Equation~\eqref{shear:eq:257-marginal} makes the summed error on any
partition into eight equal parts \(o(1)\). Exact trimming gives \(f\le\mu=q_d^{*T}\)
of mass \(z=1-o(1)\) and cap one. The orbit bound is
\(\|\widehat f(\rho)\|_{\op}\le\sqrt{8C_2}D^{-1/4}\).
In \eqref{shear:eq:subprob-plancherel}, the trivial term is
\((z^8-1)^2\), the sign term is at most \((1-z)^{16}\), and the
remaining sum is at most \((8C_2)^8Z_2(N)\). These terms and the
restored mass \(1-z^8\) all vanish. The eight time blocks take \(2056d\)
physical shuffles.

For the random omitted set route, Proposition~\ref{shear:prop:reservoir}
and the partition averaging above give a partition into eight equal parts with
summed error \(o(1)\) at \(T=257d\).
Heavy-coset deletion loses \(\eta=o(1)\) and gives cap
\(K=2/(1-\eta)\). The orbit bound is
\(\sqrt{16C_2/(1-\eta)}D^{-1/4}\). The eight-factor nonsign sum is
at most \([16C_2/(1-\eta)]^8Z_2(N)=o(1)\), the sign term is at most
\((2\eta)^{16}\), and restoring the true factors costs \(8\eta=o(1)\).
Again \(8T=2056d\). Equation~\eqref{eq:worst-state} handles every start.
\end{proof}

\begin{proposition}[Sixteen time blocks from averaged input lists]
\label{frames:average-marginals}
Let \(d\ge6\), \(b=64\), \(q=N(1-1/b)\), \(T=2048d\), and
\(\mu_T=q_d^{*T}\). The average over uniform ordered lists of \(q\)
distinct input labels of the total-variation distance of their outputs
from a uniform injection is \(o(1)\). There is a subprobability
\(f\le\mu_T\) of mass \(1-o(1)\) and a partition into 64 equal parts
for which every complementary marginal of \(f\) is bounded above by
the corresponding uniform marginal. Consequently,
\[
 \|q_d^{*(32768d)}-U_G\|_{\TV}=o(1).
\]
The full-deck conclusion holds for every deterministic starting deck.
\end{proposition}
\begin{proof}
The ordinary marginal assertion is Lemma~\ref{frames:average-law},
specifically \eqref{frames:average-tv}. The partition averaging above
selects 64 equal parts whose complementary errors sum to \(o(1)\).
Exact trimming by Proposition~\ref{shear:prop:fourier-inputs} gives
\(0\le f\le\mu_T\), of mass \(z=1-o(1)\), with cap one for all 64
coset systems. This is the stated domination of the complementary
marginals.

Use the direct inverse-tenth input cited in
Proposition~\ref{shear:prop:fourier-inputs} and the orbit estimate
\eqref{shear:eq:noncovering-orbit-general} with
\(r=64\), \(u=10\), and \(z\le1\). It gives
\[
 \|\widehat f(\rho)\|_{\op}
 \le\sqrt{64C_{10}}D^{-13/32},
 \qquad -\frac12+\frac{12}{128}=-\frac{13}{32}.
\]
At sixteen factors the nonsign, nontrivial Plancherel sum is at most
\[
 (64C_{10})^{16}
 \sum_{\substack{\lambda\vdash N\\
                  \lambda\notin\{(N),(1^N)\}}}
 D_\lambda^{\,2-32(13/32)}
 \le(64C_{10})^{16}Z_{10}(N)=o(1),
\]
because the displayed exponent is \(-11\). The sign contribution is
at most \((1-z)^{32}\), and the trivial term is \((z^{16}-1)^2\).
Equation~\eqref{shear:eq:subprob-plancherel} and restoration of the
missing mass \(1-z^{16}=o(1)\) prove convergence for \(\mu_T^{*16}\).
The block length \(T=2048d\) is a multiple of \(d\), so these are
\(16T=32768d\) physical shuffles. Equation~\eqref{eq:worst-state}
gives the assertion from every deterministic start.
\end{proof}

\begin{theorem}[Thirty-two time blocks from coefficient projection]
\label{shear:thm:32-block}
The fixed-list marginal bound at \(T=2049d\) has a separate
Hilbert--Schmidt amplification proving
\[
 \|q_d^{*(65568d)}-U_G\|_{\TV}=o(1).
\]
\end{theorem}
\begin{proof}
Use 32 equal label parts. Equation~\eqref{shear:eq:2049-marginal}
makes all complementary errors \(o(1)\). Exact trimming gives
\(f\le\mu=q_d^{*T}\) of mass \(z=1-o(1)\). Its normalization
\(\nu=f/z\) is within \(1-z\) of \(\mu\) in TV and has cap at most
two for large \(d\). The coefficient estimate
\eqref{shear:eq:hs-general}, with \(r=32\) and \(K=2\), gives
\[
 \|\widehat\nu(\rho)\|_{\HS}^2
 \le C'D^{-5/8},\qquad C'=128C_{10}.
\]
For 32 factors the nonsign, nontrivial Plancherel sum is at most
\[
 (C')^{32}\sum_{\substack{\lambda\vdash N\\
                           \lambda\notin\{(N),(1^N)\}}}D_\lambda^{-19}
 \le(C')^{32}Z_{10}(N)=o(1).
\]
The sign term is at most \([2(1-z)]^{64}\), and the trivial coefficient
is one. Restoring the true factors costs \(32(1-z)=o(1)\).
The physical time is \(32T=32(2049d)=65568d\); translation invariance
gives the same conclusion from every deterministic start.
\end{proof}

For the last application we use the companion's eventual type-count
bound. For every fixed integer \(r\ge1\), uniformly over
\(\lambda\vdash N\) with \(k_\lambda\ge1\), all sufficiently large \(N\)
satisfy
\begin{equation}
 4rM(k_\lambda)\le D_\lambda^{1/4}.
 \label{shear:eq:type-absorption}
\end{equation}
This is \cite[Lemma~\Lref{l:degree-type-absorption}]{partial-fourier}.

\begin{theorem}[Four time blocks from the branching estimate]
\label{shear:thm:4096}
The averaged-input estimate of Proposition~\ref{shear:prop:random-tuple}
has an amplification proving
\[
 \|q_d^{*(4096d)}-U_G\|_{\TV}=o(1).
\]
\end{theorem}
\begin{proof}
At \(T=1024d\), choose eight equal label parts whose summed complementary
error is \(o(1)\). Heavy-coset deletion produces a probability \(\nu\)
within \(\eta=o(1)\) of \(\mu=q_d^{*T}\), with cap at most four.
For every \(\lambda\notin\{(N),(1^N)\}\), the branching estimate and
\eqref{shear:eq:type-absorption}, with \(r=8\), give
\[
 \|\widehat\nu(\rho_\lambda)\|_{\HS}^2
 \le4rM(k_\lambda)D_\lambda^{-1+2/r}
 \le D_\lambda^{-1/2}.
\]
Four factors therefore leave the inverse-first Plancherel sum:
\[
 4\|\nu^{*4}-U_G\|_{\TV}^2
 \le(2\eta)^8+
       \sum_{\substack{\lambda\vdash N\\k_\lambda\ge1}}D_\lambda^{-1}
 =(2\eta)^8+Z_1(N)=o(1).
\]
The displayed sign term uses \eqref{shear:eq:parity}; the trivial
coefficient is one. Restoring the four true factors costs \(4\eta=o(1)\).
Their physical time is \(4T=4096d\), and \eqref{eq:worst-state}
gives the same conclusion from every start.
\end{proof}

\begin{corollary}[Mixing order and lower estimates]
\label{shear:cor:order}
For all sufficiently large \(d\),
\[
 d<t_{\mathrm{mix}}(d)\le512d.
\]
Thus the mixing time has order \(\Theta(d)\) in physical shuffles.
The one-card support bound also gives \(t_{\mathrm{mix}}(d)\ge d\)
for every \(d\ge1\).
\end{corollary}
\begin{proof}
Lemma~\ref{lem:lower} gives distance at least
\(1-(e/\sqrt N)^N>1/4\) for every integer \(t\le d\) when \(d\) is
large. Lemma~\ref{lem:one-card-lower} gives the stated one-card bound
for every \(d\). The upper bound follows from
Corollary~\ref{shear:cor:512}; \eqref{eq:sharp-support-time} gives
the sharper lower asymptotic \(2d-O(1)\).
\end{proof}
